\section{\uppercase{Introduction}} \label{sec1}
Large Language Models (LLMs) perform well across a wide variety of tasks~\cite{Zhao25}. However, because LLM outputs rely heavily on knowledge in their training data, they can generate misinformation or hallucinate when queried about unlearned knowledge~\cite{Ji23}. Furthermore, LLMs struggle to learn continuously updated real-world information~\cite{Wang25}. To address these challenges, Retrieval-Augmented Generation (RAG)~\cite{Lewis20} has been proposed, which retrieves external knowledge and presents it as context to the LLM. RAG enables LLMs to generate responses that reflect up-to-date, accurate knowledge.

One representative application of RAG is question answering (QA) tasks. Early research on RAG-based QA focused primarily on single-hop QA tasks~\cite{rajpurkar16}, where the answer exists within a single document~\cite{Shi24}. However, many multi-hop QA tasks require cross-document association and step-by-step reasoning~\cite{yang18,ho20,trivedi22}. For example, the query ``Who was the U.S. President when Einstein published the theory of relativity?'' requires a two-step reasoning process: identifying the publication year and matching it with the serving President. To handle such complex questions, researchers have proposed specialized RAG methods for multi-hop QA tasks, such as LightRAG~\cite{Guo25} and GraphRAG~\cite{edge25}.

However, these specialized methods suffer from increased inference latency because they require multiple retrieval steps and multi-step reasoning. In real-world operations, not all input queries are complex. Many frequently asked questions are simple and direct; for instance, a question such as ``When is Einstein's birthday?'' can often be answered solely using the LLM's parametric internal knowledge without retrieving external documents~\cite{Mallen23}. For such simple queries, unnecessary retrieval and reasoning steps merely incur extra computational cost. However, processing a complex question with an overly simplistic RAG configuration may result in a degradation of response quality and accuracy~\cite{LingjiaoChen24}.

% Research Challenges
Against this background, research on ``adaptive RAG'' that dynamically adjusts RAG components according to query characteristics has advanced rapidly~\cite{Asai23,fu24,Tang25,Rishi25,zhang2025adagcrag}. However, existing adaptive RAG approaches face three main challenges:
First, agentic selection methods (such as Agentic RAG)~\cite{Asai23} that rely on LLM self-judgment or planning require multi-step reasoning or online exploration to decide configurations. Although this can improve response accuracy, it adds substantial LLM inference overhead (latency and API cost) at each decision step, making it difficult to adjust multiple micro-level hyperparameters (e.g., chunk size and overlap) comprehensively and cost-effectively.
Second, methods using online reinforcement learning~\cite{fu24,Tang25} involve continuous trial-and-error in live environments, increasing operational infrastructure costs and learning instability.
Third, while adaptive approaches based on query or document features exist~\cite{Rishi25,zhang2025adagcrag}, most focus on switching macro-level processing paradigms or controlling parameters of specific graph-based RAGs. Prior work has not adequately addressed unified ranking over a large discrete configuration space (96 combinations) spanning heterogeneous RAG paradigms (Graph/Non-Graph) and micro-level hyperparameters such as chunk size and overlap. Furthermore, frameworks that dynamically control opposing multi-objective metrics (precision, recall, $F_1$ score, and execution latency) during inference according to user-defined priority weights $W$ remain scarce.

% Research Objectives
To address these challenges, this study adopts an offline learning approach that learns optimal configurations from pre-collected empirical data, eliminating reliance on LLMs or continuous online exploration. Instead of runtime adjustments that require extra LLM inference, we train ranking models (RM) to optimize both accuracy and cost over a predefined candidate RAG configuration space. At inference time, the trained RM instantly ranks candidate configurations without LLM overhead, and we select the optimal configuration based on user-defined priority weights. Compared with Agentic RAG involving online search, our method significantly reduces inference cost and decision latency (millisecond-order selection) while accommodating a broad configuration space encompassing multiple RAG types and hyperparameters. Moreover, because it does not rely on online learning, it avoids ongoing operational overhead. Furthermore, by incorporating surface-level and semantic features from both queries and documents to learn multi-objective metrics, our approach effectively addresses the third challenge.

% Overview of Methodology
In the proposed method, we first extract statistical surface-level features and embedding-based semantic features from the query and document. During training, we train the RM using empirically measured performance across diverse RAG configurations as ground-truth data. At inference time, the RM ranks candidate configurations for each input query, allowing the system to select cost-effective or high-accuracy RAG configurations. In addition, users can dynamically adjust priority weights according to operational demands (e.g., accuracy-oriented, balanced, or cost-oriented), enabling flexible multi-objective decision-making.

% Overview of Experiments
In the evaluation experiments, we assess the RM's ranking performance on QA tasks and the final response quality generated by RAG under the selected configurations.

The contributions of this study are as follows:
\begin{itemize}
\item We experimentally show that while RAG configuration information (e.g., generative models and retrieval methods) serves as the primary basis for ranking, integrating surface-level and semantic features of both queries and documents achieves the highest ranking accuracy (NDCG@1=68.52\%, Hit@5=65.00\%), outperforming configuration information alone (NDCG@1=66.84\%).
\item We introduce a multi-objective ranking approach that unifies evaluation of accuracy (precision, recall, $F_1$ score) and cost (execution time), enabling adaptive, automatic RAG configuration selection tailored to user priority weights and input context.
\end{itemize}

The remainder of this paper is organized as follows: Section~\ref{sec2} reviews related work, Section~\ref{sec3} details the proposed method, Section~\ref{sec4} describes the experimental evaluation methodology, Section~\ref{sec5} presents the experimental results and discussion, and Section~\ref{sec6} concludes the paper and outlines future work.

%Note that this paper extends the preliminary study by Anaguchi et al.~\cite{anaguchi26} with the following specific additions and enhancements: First, we clarify the positioning of our work by comparing it with recent related systems (CARROT, uRAG, RAGRouter, and AdaGCRAG). Second, we incorporate Adaptive-RAG, a representative adaptive RAG baseline, into our evaluation framework and evaluate performance using Pareto frontiers between accuracy and cost. Third, in addition to automatic evaluation metrics (BERTScore), we re-evaluate the proposed method using standard QA metrics, specifically Exact Match (EM) and Token-level $F_1$ score. Fourth, we perform Wilcoxon signed-rank tests to confirm the statistical significance of performance differences between our method and baselines. Finally, we quantify the sequential execution cost in top-$K$ candidate presentation models as cumulative execution time, demonstrating practical cost efficiency.